% =====================================================================
%  4 MÉTODO DE PESQUISA — "Como resolvi o problema?"
%  Dois tópicos (Item 3): (1) classificação da pesquisa;
%  (2) estruturação, SEMPRE representada por uma Figura (Dica 5).
%  As etapas daqui viram as subseções dos Resultados (Dica 6).
% =====================================================================
\section{Método de Pesquisa}\label{sec:metodo}

\escrever{parágrafo de ligação com o Referencial Teórico.}

\subsection{Classificação da pesquisa}

\escrever{natureza (aplicada), abordagem (quantitativa), objetivos (exploratória e experimental) e procedimentos (pesquisa experimental em laboratório). Descrever amostras: número de sensores fabricados, número de ensaios, número de voluntários (se houver), instrumentos de medição e técnicas de análise.}

\dica{Se houver testes com voluntários, citar a aprovação do \ac{CEP} (Resolução CNS 466/2012) e colocar o termo de consentimento no Apêndice.}

\subsection{Estruturação da pesquisa}

A pesquisa foi estruturada em sete etapas, conforme a Figura~\ref{fig:estrutura}.
\escrever{ajustar o texto e a figura às etapas reais.}

\begin{figure}[H]
  \centering
  \begin{spacing}{1.0}
  \begin{tikzpicture}[
      node distance=4mm and 6mm,
      etapa/.style={draw, rounded corners=2pt, align=center, font=\footnotesize,
                    minimum width=3.6cm, minimum height=1.25cm, fill=gray!6},
      seta/.style={-{Stealth[length=2mm]}, thick}]
    \node[etapa] (e1) {\textbf{Etapa 1}\\Levantamento do\\estado da arte (TEMAC)};
    \node[etapa, right=of e1] (e2) {\textbf{Etapa 2}\\Projeto CAD e\\fabricação em TPU};
    \node[etapa, right=of e2] (e3) {\textbf{Etapa 3}\\Instrumentação\\e firmware};
    \node[etapa, below=of e3] (e4) {\textbf{Etapa 4}\\Caracterização\\em bancada};
    \node[etapa, left=of e4] (e5) {\textbf{Etapa 5}\\Calibração\\(regressão $\times$ MLP)};
    \node[etapa, left=of e5] (e6) {\textbf{Etapa 6}\\Integração na luva e\\feedback vibrotátil};
    \node[etapa, below=of e6] (e7) {\textbf{Etapa 7}\\Validação com\\objetos reais};
    \draw[seta] (e1) -- (e2);
    \draw[seta] (e2) -- (e3);
    \draw[seta] (e3) -- (e4);
    \draw[seta] (e4) -- (e5);
    \draw[seta] (e5) -- (e6);
    \draw[seta] (e6) -- (e7);
  \end{tikzpicture}
  \end{spacing}
  \caption{Estruturação da pesquisa}
  \label{fig:estrutura}
  \fonte{Elaborado pelo autor (2026).}
\end{figure}

\subsubsection{Etapa 1 -- Levantamento do estado da arte}
\escrever{como a revisão TEMAC foi conduzida (detalhes na Seção~\ref{sec:estado-arte}).}

\subsubsection{Etapa 2 -- Projeto CAD e fabricação}
\escrever{dimensões (15$\times$15$\times$10~mm), tamanho de célula, espessura das paredes, critério de imprimibilidade (bico de 0,4~mm), filamento \ac{TPU} 95A, ímãs N52, pausa de impressão para inserção dos ímãs.}

\subsubsection{Etapa 3 -- Instrumentação e firmware}
\escrever{TLV493D na lateral da falange distal, multiplexador TCA9548A, \ac{ESP32}, taxa de amostragem, linha de base ($\Delta B$), arquitetura modular do firmware.}

\subsubsection{Etapa 4 -- Caracterização em bancada}
\escrever{instrumento de referência (célula de carga), faixas de força normal (0--15~N) e cisalhante (0--8~N), malha de localização 10$\times$10~mm, número de repetições, controle de temperatura.}

\subsubsection{Etapa 5 -- Calibração}
\escrever{modelos comparados, divisão dos dados por ensaio (validação temporal), métricas, exportação do modelo para o microcontrolador.}

\subsubsection{Etapa 6 -- Integração e feedback vibrotátil}
\escrever{posição dos motores \ac{ERM} (dorso ou palma), driver MOSFET, mapeamento força $\rightarrow$ \ac{PWM}, teste de interferência magnética em função da distância, medição de latência.}

\subsubsection{Etapa 7 -- Validação}
\escrever{tarefas (aperto controlado, detecção de deslizamento e pinça fina), condição de comparação, métricas e critério de sucesso.}
